\section{Sharp finite-experiment budgets}\label{sec:v7-finite-experiments}
For experiments $\mathcal E=\{P_\theta\}$ and $\mathcal F=\{Q_\theta\}$
with the same parameter set, write
\[
 \delta(\mathcal E,\mathcal F)=\inf_K\sup_\theta
             \|KP_\theta-Q_\theta\|_{\rm TV},\qquad
 \Del(\mathcal E,\mathcal F)=
             \max\{\delta(\mathcal E,\mathcal F),
                    \delta(\mathcal F,\mathcal E)\}.
\]
All parameterized laws in this section are Borel families.
The infimum is over parameter-independent Markov kernels. These are
terminal experiment comparisons in the sense of
\cite{LeCamYang2000,Torgersen1991}; only the encoding statistic, not the
comparison randomization, is part of the physical data reduction.

Let $\mathcal G=\{N(z,\Gamma):z\in K\}$, where
$K\subset\R^r$ is compact and $\Gamma$ is positive definite. In the
lower bound assume that $K$ contains a nonempty Euclidean ball.

\begin{theorem}[Gaussian finite-experiment cardinality]
\label{thm:v7-gaussian-budget}
There are constants $c,C>0$, depending on $K,\Gamma,r$, such that
\begin{equation}
 cM^{-2/r}\le
 \inf_{\mathcal F:\,|\mathcal F|\le M}\Del(\mathcal F,\mathcal G)
 \le CM^{-2/r},\qquad M\ge1.                            \label{eq:v7-gaussian-budget}
\end{equation}
The lower bound ranges over all experiments on at most $M$ outcomes,
not just deterministic Gaussian quantizers. The upper bound is attained
in order by a deterministic statistic with $J^r$ labels, where
$J=\lfloor M^{1/r}\rfloor$, followed by a specified positive
reconstruction kernel. No overflow label is needed.
\end{theorem}

\subsection{The cardinality obstruction}
\begin{lemma}[A decision-theoretic lower bound]\label{lem:v7-cardinality}
Under the preceding assumptions,
$\Del(\mathcal F,\mathcal G)\ge cM^{-2/r}$ whenever
$\mathcal F$ has at most $M$ outcomes.
\end{lemma}
\begin{proof}
Restrict the mean parameter to a closed cube $K_0$ of positive side
length contained in $K$ and put the uniform prior on $K_0$. Restriction
can only decrease an experiment distance. Use squared-error estimation
of $z$, with actions restricted to $K_0$; its loss range is bounded by
$L=\diam(K_0)^2$. Let $m(y)=\E[z\mid y]$ in this Gaussian experiment.
Differentiation under the compact prior integral gives
\begin{equation}
                    Dm(y)=\Cov(z\mid y)\Gamma^{-1}.    \label{eq:v7-posterior-jacobian}
\end{equation}
The posterior density is positive on the cube. Hence its covariance is
positive definite, and $Dm(y)$ is nonsingular for every $y$. The inverse
function theorem on a small neighborhood of any fixed $y_0$, together
with the positive continuous marginal density of $y$, shows that the
law of $m(y)$ has a density bounded below by a positive constant on
some smaller cube in $\R^r$.

This implies $e_M^2(m)\ge c_0 M^{-2/r}$. To see it directly, put balls
of radius $t=(|Q|/(2M v_r))^{1/r}$ about the $M$ proposed centres, where
$Q$ is the smaller cube and $v_r$ the volume of the unit ball. At least
half of $Q$ remains outside them, and contributes at least a fixed
multiple of $t^2$. This argument permits arbitrary centres.

Every $M$-outcome garbling $S$ of the Gaussian observation therefore has
Bayes risk at least $B_{\mathcal G}+c_0M^{-2/r}$, by conditional
orthogonality. If $\Del(\mathcal F,\mathcal G)<\epsilon$, choose a
kernel $K$ from $\mathcal G$ to $\mathcal F$ with error below $\epsilon$.
The experiment $K\mathcal G$ has at most $M$ outputs, and bounded-loss
comparison gives
\[
 B_{K\mathcal G}\le B_{\mathcal F}+L\epsilon
                  \le B_{\mathcal G}+2L\epsilon.
\]
The second inequality uses a kernel in the reverse direction. Comparing
these bounds and taking the infimum over $\epsilon$ proves the result.
Randomized actions do not improve a Bayes squared-error minimum. Notice
that $\mathcal F$ itself was not assumed to be a Gaussian garbling.
\end{proof}

\subsection{Positive second-order reconstruction}
Put
\begin{equation}
 \tau(x)=\tfrac12+\pi^{-1}\arctan x,
 \qquad u=\tau^{\otimes r}(y)\in(0,1)^r.               \label{eq:v7-compactification}
\end{equation}
Let $p_z$ be the transformed Gaussian density. If $x_j=\tan\pi(u_j-1/2)$,
then
\begin{equation}
 p_z(u)=\phi_\Gamma(x-z)\prod_{j=1}^r\pi(1+x_j^2).     \label{eq:v7-transformed-density}
\end{equation}
It extends by zero across every face to a smooth periodic density.
Indeed each fixed-order derivative is a polynomial in $x$ times the
Gaussian density, which tends to zero faster than any polynomial as a
face is approached. Uniformity holds for $z$ in the fixed compact set.
For use below define the finite constants
\begin{equation}
 H_2=\sup_{z,u}\|D^2p_z(u)\|_{\rm op},\qquad
 H_0=\max_j\sup_{z,u_j}p_{z,j}(u_j),                    \label{eq:v7-Hconstants}
\end{equation}
where $p_{z,j}$ is a transformed marginal density. These suprema, or
any upper bounds obtained from \eqref{eq:v7-transformed-density}, fix
all the reconstruction and boundary constants; they do not depend on
$M$ or sample size. For instance, differentiating twice introduces only
polynomials and the eigenvalue bounds of $\Gamma$, so bounded mean
norm and fixed covariance conditioning suffice.

Partition the flat unit torus into $J^r$ cubes of side $h=1/J$, with
centres $c_i$, and write $Q_J(u)$ for the cube label. For $J\ge2$ let
$\Phi_i$ be the periodic tensor product of the piecewise-linear hat
functions with value one at $c_i$, zero at its neighboring mesh centres,
and support within one mesh step in each coordinate. Then
\[
 \Phi_i\ge0,\qquad \int\Phi_i=h^r,\qquad
 \sum_i\Phi_i(u)=1.
\]
On each mesh interval the neighbouring centres reproduce the affine
coordinate, using local periodic lifts across the seam. Define the
reconstruction from label $i$ to have density $h^{-r}\Phi_i(u)$,
then apply $\tau^{-1}$ coordinatewise. Values on the faces have
probability zero.

\begin{lemma}[A quadratic reconstruction estimate]
\label{lem:v7-hat}
Uniformly in $z\in K$, this kernel $H_J$ satisfies
\begin{equation}
 \|H_J Q_J\tau^{\otimes r} N(z,\Gamma)-N(z,\Gamma)\|_{\rm TV}
                         \le\frac{r H_2}{12J^2}.       \label{eq:v7-hat}
\end{equation}
\end{lemma}
\begin{proof}
Let $p_i$ be the probability of cube $i$. Symmetric Taylor expansion
about its centre gives
$|p_i/h^r-p_z(c_i)|\le rH_2h^2/24$.
The linear interpolant $I_hp_z=\sum_i p_z(c_i)\Phi_i$ has error at most
$rH_2h^2/8$: its weights reproduce constants and affine functions and
the sum of their coordinate variances is at most $rh^2/4$. The second
Taylor remainder is half this variance times $H_2$. Both calculations
hold across the seam because $p_z$ has a smooth periodic extension.
The reconstructed density is
$\sum_i(p_i/h^r)\Phi_i$. Nonnegativity and the partition of unity give
supremum error at most $rH_2h^2/6$. Integrating over volume one and
using the factor $1/2$ for total variation proves the claim. Inverting
\eqref{eq:v7-compactification} preserves this distance.
\end{proof}

\begin{proof}[Proof of Theorem~\ref{thm:v7-gaussian-budget}]
Lemma~\ref{lem:v7-cardinality} gives the lower bound. The deterministic
statistic $Q_J\tau^{\otimes r}$ has $J^r$ outputs, so the forward
deficiency is zero; Lemma~\ref{lem:v7-hat} gives the reverse deficiency.
For $M\ge2^r$, $J=\lfloor M^{1/r}\rfloor\ge M^{1/r}/2$.
For the remaining finite range use any one-outcome experiment and the
trivial bound $\Del\le1$, increasing $C$. The reconstruction uses
positive probabilities; it is not a signed interpolation operator.
\end{proof}

The use of a compactified grid avoids a growing Gaussian overflow
radius. The use of hats rather than constant cells cancels the first
Taylor term. Both choices are necessary for the stated proof of a
logarithm-free quadratic budget rate.
